% !PRINTED 67
\weberpage{67}
%<<CONT

\gesperrt{III. Die Permutationen der ersten Art sind aus einer geraden Anzahl von Transpositionen zusammengesetzt, und die der zweiten Art aus einer ungeraden Anzahl.}

Zu der ersten Art ist dann auch die sogenannte identische Permutation zu rechnen, die \(\mathfrak A\) ungeändert lässt.

Daraus ergiebt sich noch die Folgerung: Auf wie verschiedenen Wegen man auch \(\mathfrak A'\) aus \(\mathfrak A\) durch Transpositionen ableiten mag, die Anzahl dieser Transpositionen ist bei allen diesen Arten übereinstimmend gerade oder ungerade (je nachdem \(\mathfrak A'\) zur ersten oder zur zweiten Art gehört).

Wenn wir in den sämmtlichen Anordnungen
\[
\mathfrak A,\mathfrak A',\mathfrak A'',\ldots
\]
der \(n\) Elemente eine Transposition, etwa \((1,2)\), vornehmen, so geht jede dieser Anordnungen in eine bestimmte andere über, etwa \(\mathfrak A\) in \(\mathfrak B\), \(\mathfrak A'\) in \(\mathfrak B'\), \(\mathfrak A''\) in \(\mathfrak B''\) \ldots, und wenn wir dieselbe Transposition noch einmal wiederholen, so geht \(\mathfrak B\) wieder in \(\mathfrak A\), \(\mathfrak B'\) wieder in \(\mathfrak A'\) \ldots\ über. Daraus folgt, dass die Anordnungen \(\mathfrak B,\mathfrak B',\mathfrak B'',\ldots\) alle von einander verschieden sind und folglich in ihrer Gesammtheit mit der Gesammtheit der \(\mathfrak A\) übereinstimmen. Da nun, wie wir oben gesehen haben, die sämmtlichen Differenzenproducte
\[
P,P',P'',\ldots
\]
die aus \(P\) mit den verschiedenen Anordnungen \(\mathfrak A,\mathfrak A',\mathfrak A'',\ldots\) gebildet sind, durch eine Transposition das Zeichen ändern, so folgt, dass jedem \(\mathfrak A\) der ersten Art ein \(\mathfrak B\) der zweiten Art entspricht und jedem \(\mathfrak A\) der zweiten Art ein \(\mathfrak B\) der ersten Art.

\gesperrt{IV. Hiernach ist die Anzahl der Anordnungen der ersten Art ebenso gross, wie die Anzahl der Anordnungen der zweiten Art, nämlich} \({1\over2}\Pi(n)\)\footnote{Diese Sätze sind hier aus der Betrachtung des Productes \(P\), also einer Zahlgrösse, gewonnen. Wie man ohne Benutzung einer solchen Function zu denselben Ergebnissen gelangen kann, werden wir im XIV. Abschnitt sehen.}.

Für \(n=3\) haben wir die folgenden sechs Anordnungen, von denen die erste Horizontalreihe die erste Art bildet:
\begin{equation}
\begin{array}{ccc}
(1,2,3),&(2,3,1),&(3,1,2),\\
(3,2,1),&(2,1,3),&(1,3,2).
\end{array}
\tag{3}
\end{equation}
